\documentclass[12pt,a4paper]{report}
\usepackage[utf8]{inputenc}
\usepackage[spanish]{babel}
\usepackage{geometry}
\geometry{margin=2.5cm}
\usepackage{setspace}
\onehalfspacing
\usepackage{hyperref}

\title{\textbf{Biblia del Universo}\\[0.5cm]\large El Velo Gravitacional}
\author{Carlos G. Ramirez}
\date{31 de diciembre de 2025}

\begin{document}

\maketitle
\tableofcontents
\newpage

\chapter{BIBLIA DEL UNIVERSO}
\section\textit{{El Velo Gravitacional}

> \textbf{Documento de Referencia Canónico}  
> Última actualización: 29 de diciembre de 2025 (v1.3)  
> Autor: [Tu nombre]

---

\section}{📋 ÍNDICE}

1. [Sinopsis General](#sinopsis-general)
2. [Personajes Principales](#personajes-principales)
3. [Personajes Secundarios](#personajes-secundarios)
4. [Lugares y Geografía](#lugares-y-geografía)
5. [Tecnología](#tecnología)
6. [Cultura y Tradiciones](#cultura-y-tradiciones)
7. [Historia y Cronología](#historia-y-cronología)
8. [Eventos Clave](#eventos-clave)
9. [Temas y Motivos](#temas-y-motivos)

---

\section\textit{{SINOPSIS GENERAL}

\textbf{El Velo Gravitacional} es una novela de ciencia ficción distópica ambientada en el año 2114, en un mundo devastado por el cambio climático antropogénico. La historia sigue a Ethan Ramírez, un joven piloto de 17 años del marginalizado Distrito RD, quien compite en una carrera mortal de manipulación gravitacional para ganar la ciudadanía en la privilegiada Goglee City. La narrativa explora temas de desigualdad social, identidad cultural, resistencia y el costo humano del progreso tecnológico descontrolado.

---

\section}{PERSONAJES PRINCIPALES}

\subsection\textit{{ETHAN RAMÍREZ}
\begin{itemize}
    \item \textbf{Físico:} 17 años, piel clara de tono cálido, rasgos con un sutil toque latino. Cabello rojo-fresa/cobrizo, corto y ligeramente crespo. Complexión atlética y robusta, desarrollada por años de trabajo físico en el taller.

  \item \textbf{Fecha de nacimiento}: 3 de febrero de 2095

  \item \textbf{Lugar de nacimiento}: Distrito RD, periferia de Goglee City

  \item \textbf{Descripción física}: [A desarrollar en la narrativa]

  \item \textbf{Personalidad}: 
\end{itemize}
  - Determinado y soñador
  - Talento natural para la mecánica
  - \textbf{Habilidad especial}: Puede "sentir" la gravedad antes de que cambie
  - \textbf{Don único}: Tiene la capacidad de "ver" las geodésicas (caminos curvos en el espacio-tiempo)
    - Este don es extremadamente raro (solo 2 personas conocidas en 30 años: Luna Reyes y Ethan)
    - Se manifiesta como una percepción extrasensorial de la geometría del espacio-tiempo
    - \textbf{Visiones del Velo}: Sufre de sueños confusos y fragmentados que parecen visiones en primera persona desde una cabina de mando avanzada. Ve interfaces HUD de tipo casco futurista y distorsiones espaciales en forma de túneles o agujeros de gusano. Estos sueños son un misterio recurrente; solo cobran sentido una vez que Ethan se enfrenta a la realidad de la carrera o del cataclismo. Debe mencionar estas experiencias al menos dos veces en la historia para crear intriga.
  - Leal a sus raíces pero anhela un futuro mejor.
\begin{itemize}
  \item \textbf{Historia personal}:
\end{itemize}
  - Madre desapareció cuando tenía ~5 años (2100)
  - Criado solo por su padre Carlos
  - Recibió su primera motowave como regalo de cumpleaños (17 años)
  - Pasó 10 meses transformándola en "El Fénix"
\begin{itemize}
  \item \textbf{Objetivo}: Ganar El Velo Gravitacional para conseguir ciudadanía y mejorar la vida de su padre y su gente

  \item \textbf{Arco narrativo}: De soñador ingenuo a líder revolucionario
\end{itemize}

\subsection}{CARLOS RAMÍREZ}
\begin{itemize}
  \item \textbf{Edad}: 56 años (al inicio de la historia principal)

  \item \textbf{Fecha de nacimiento}: 3 de enero de 2058 (año de La Gran Inundación)

  \item \textbf{Lugar de nacimiento}: Distrito RD (hijo de refugiados caribeños)

  \item \textbf{Descripción física}: 
\end{itemize}
  - Robusto, piel oscura profunda.
  - Pelo ensortijado gris/canoso, barba densa con canas marcadas.
  - Gafas de pasta negra.
  - Vestimenta: Estilo ciberpunk/steampunk funcional (camisas de cuello reforzado, elementos mecánicos).
  - Manos marcadas por años de trabajo mecánico.
\begin{itemize}
  \item \textbf{Personalidad}:
\end{itemize}
  - Protector y sabio
  - Ateo recalcitrante (aunque mantiene el lenguaje cultural cristiano)
  - Pragmático y basado en la evidencia física
  - Maestro paciente
  - Lleva el peso de la historia de su pueblo
\begin{itemize}
  \item \textbf{Profesión}: Mecánico experto y maestro del Distrito RD

  \item \textbf{Historia personal}:
\end{itemize}
  - Nacido el mismo año de La Gran Inundación
  - Nunca conoció la isla de sus padres
  - Tuvo a Ethan a los 37 años
  - Esposa desapareció cuando Ethan tenía ~5 años
  - Crió solo a su hijo
\begin{itemize}
  \item \textbf{Relación con Ethan}: Padre amoroso que le enseñó mecánica, cultura y valores

  \item \textbf{Herencia cultural}: Descendiente de refugiados del Caribe (isla específica sin confirmar)
\end{itemize}

---

\section\textit{{PERSONAJES SECUNDARIOS}

\subsection}{CÍRCULO CERCANO DE ETHAN (PERIODO DE APRENDIZAJE)}

#### YEISLY RAMÍREZ
\begin{itemize}
  \item \textbf{Relación}: Prima de Ethan (hija de hermano/hermana de Carlos)

  \item \textbf{Edad}: 18 años (1 año mayor que Ethan)

  \item \textbf{Fecha de nacimiento}: ~2096

  \item \textbf{Descripción física}: [A desarrollar]

  \item \textbf{Personalidad}: 
\end{itemize}
  - Artista urbana talentosa
  - Protectora como hermana mayor
  - Creativa y expresiva
  - Voz de la resistencia cultural
\begin{itemize}
  \item \textbf{Habilidad}: Graffiti y arte urbano

  \item \textbf{Rol en la historia}: 
\end{itemize}
  - Hermana mayor que Ethan nunca tuvo (especialmente importante dado que la madre desapareció)
  - Documenta la construcción de El Fénix en murales
  - Sus murales se convierten en símbolos de resistencia
  - Conexión familiar fuerte con Ethan y Carlos
  - Apoya emocionalmente a ambos
\begin{itemize}
  \item \textbf{Relación con Ethan}: Cercana, se cuentan todo, ella lo anima cuando duda

  \item \textbf{Subplot potencial}: Su arte puede meterla en problemas con las autoridades corporativas
\end{itemize}

#### RAFA (RAFAEL)
\begin{itemize}
  \item \textbf{Edad}: 17 años

  \item \textbf{Descripción física}: [A desarrollar]

  \item \textbf{Personalidad}: 
\end{itemize}
  - Leal hasta la muerte
  - Optimista y bromista
  - Equilibra la seriedad de Ethan
  - Más social, saca a Ethan de su zona de confort
\begin{itemize}
  \item \textbf{Habilidad}: Experto en electrónica y sistemas

  \item \textbf{Rol en la historia}: 
\end{itemize}
  - Mejor amigo de Ethan
  - Ayuda con los circuitos y sistemas eléctricos de El Fénix
  - Apoyo emocional constante
  - Quizás también quiere competir pero no tiene recursos
  - Vive el sueño a través de Ethan
\begin{itemize}
  \item \textbf{Relación con Ethan}: Hermanos de corazón, se conocen desde niños. Rafa es la sombra técnica de Ethan; siempre está cerca cuando hay que ensuciarse las manos con "El Fénix".

  \item \textbf{Lenguaje}: Usa constantemente slang caribeño activo. Su saludo de cabecera es "klk", "klk manito" o "klk manín". Es la voz que mantiene el ánimo arriba con su energía dominicana.
\end{itemize}

#### TOMÁS "EL VIEJO TORO" VEGA
\begin{itemize}
  \item \textbf{Edad}: 68 años

  \item \textbf{Fecha de nacimiento}: ~2046

  \item \textbf{Descripción física}: 
\end{itemize}
  - Marcas de quemaduras y cicatrices de accidentes
  - Cuerpo deteriorado pero aún fuerte
  - Ojos que han visto demasiado
  - Manos temblorosas (trauma post-accidente)
\begin{itemize}
  \item \textbf{Personalidad}: 
\end{itemize}
  - Amargado pero sabio
  - Protector de Ethan (como el hijo que nunca tuvo)
  - Lleva el peso de sus fracasos
  - Mezcla de esperanza y cinismo
  - Atormentado por el pasado
\begin{itemize}
  \item \textbf{Historia}: 
\end{itemize}
  - Leyenda del Distrito RD
  - Compitió en las primeras ediciones de El Velo (2078-2085)
  - \textbf{Ganó El Velo en 2082} pero \textbf{rechazó la ciudadanía}
  - Fue castigado/exiliado por las corporaciones
  - Perdió todo: familia, hogar, respeto
  - \textbf{Tragedia de Luna Reyes (2084)}: 
    - Su mejor amiga y compañera de carrera
    - Luna murió usando El Surfeo de Geodésicas
    - Error de cálculo de 3 milímetros causó Efecto Honda
    - Tomás escuchó sus gritos por radio durante 30 segundos
    - Vio su motowave desintegrarse
    - Nunca se perdonó no haberla detenido
    - Dejó de competir después de su muerte
  - Ahora vive en las ruinas del Distrito RD
\begin{itemize}
  \item \textbf{Habilidad}: Maestro piloto, conoce todos los trucos del Velo (incluyendo El Surfeo de Geodésicas)

  \item \textbf{Rol en la historia}: 
\end{itemize}
  - Mentor técnico y emocional de Ethan
  - Le enseña trucos de pilotaje avanzados
  - Le advierte sobre los peligros del Velo
  - Le enseña a "sentir" la geometría del espacio-tiempo
  - \textbf{Le enseña El Surfeo de Geodésicas} (con gran conflicto interno)
  - Ve en Ethan una segunda oportunidad
  - Precedente histórico del rechazo de Ethan
\begin{itemize}
  \item \textbf{Relación con Ethan}: Figura paterna secundaria, quiere a Ethan como al hijo que nunca tuvo

  \item \textbf{Frases características}: 
\end{itemize}
  - "La gravedad no perdona, mijo. Pero tampoco olvida."
  - "Luna decía que era como escuchar una canción que nadie más podía oír."
\begin{itemize}
  \item \textbf{Conflicto interno}: 
\end{itemize}
  - Tiene miedo de que Ethan repita los errores de Luna
  - Pero sabe que sin El Anclaje Fantasma, Ethan no puede ganar
  - Debe elegir entre proteger a Ethan o darle la única herramienta para triunfar

#### UNIKA
\begin{itemize}
  \item \textbf{Edad}: 17 años

  \item \textbf{Descripción física}: [A desarrollar]

  \item \textbf{Personalidad}: 
\end{itemize}
  - Mecánica talentosa y de mente rápida
  - Competitiva
  - Directa y sin filtros (estilo educado y técnico)
  - Guarda sus sentimientos profundos tras una fachada de dureza
\begin{itemize}
  \item \textbf{Historia personal}:
\end{itemize}
  - \textbf{Origen}: Nacida en el Domo de Goglee City (Nivel 4).
  - Sus padres eran ingenieros de sistemas corporatvios que fueron "deseleccionados" tras denunciar corrupción.
  - Perdió la ciudadanía y fue enviada al Distrito RD con lo puesto.
  - Dedica su genio técnico a la venganza contra el sistema que la traicionó.
  - Desde pequeña ha admirado a Ethan por su valentía.
\begin{itemize}
  \item \textbf{Rol en la historia}: 
\end{itemize}
  - Ingeniera de sistemas (especialista en el Inversor G)
  - \textbf{Curiosidad Astronómica}: Su hobby es analizar archivos de data antigua de la antigua agencia espacial. No es una revelación heroica, sino una curiosidad técnica persistente.
  - \textbf{El Patrón de Ruido}: Analizando registros de la época en que la madre de Ethan desapareció (cerca del 2100), notó un "bamboleo" en una estrella lejana que fue clasificado como ruido. Recientemente, ha encontrado el mismo patrón de "ruido rítmico" pero en datos mucho más cercanos y masivos. Ella sospecha que no es estática, sino una firma gravitacional real.
  - Optimiza "El Fénix" con tecnología de grado corporativo modificada.
  - Amiga cercana de Yeisly.
\begin{itemize}
  \item \textbf{Relación con Ethan}: Amistad técnica profunda. Ella guarda sentimientos románticos secretos. Es su apoyo intelectual y técnico fundamental.
\end{itemize}

#### MARCOS
\begin{itemize}
  \item \textbf{Edad}: 17-18 años

  \item \textbf{Descripción física}: [A desarrollar]

  \item \textbf{Personalidad}: 
\end{itemize}
  - Escéptico y realista
  - Voz de la razón (a veces pesimista)
  - Leal a pesar de sus dudas
\begin{itemize}
  \item \textbf{Rol en la historia}: 
\end{itemize}
  - Parte del grupo de amigos
  - Piensa que El Velo es trampa corporativa
  - Representa la voz del escepticismo
  - Eventualmente apoya a Ethan a pesar de sus dudas
\begin{itemize}
  \item \textbf{Relación con Ethan}: Amigo que cuestiona pero apoya
\end{itemize}

---

\subsection\textit{{ANTAGONISTAS Y OTROS PERSONAJES}

#### SONY
\begin{itemize}
  \item \textbf{Edad}: 19-20 años (2-3 años mayor que Ethan)

  \item \textbf{Fecha de nacimiento}: ~2094-2095

  \item \textbf{Herencia cultural}: Descendiente de refugiados cubanos o venezolanos (costas caribeñas continentales)

  \item \textbf{Descripción física}: 
\end{itemize}
  - Pelo de colores neón (cambia frecuentemente)
  - Chaqueta de cuero con parches de bandas anarquistas
  - Mirada intensa y apasionada
  - [Más detalles a desarrollar]
\begin{itemize}
  \item \textbf{Personalidad}: 
\end{itemize}
  - Anarquista/revolucionaria apasionada
  - Experta en explosivos y sabotaje
  - Calculadora y estratégica
  - Idealista radical
  - No quiere subir al domo, quiere derribarlo
  - Protectora de Ethan pero también lo desafía
\begin{itemize}
  \item \textbf{Relación con Ethan}: 
\end{itemize}
  - \textbf{NOVIA de Ethan}
  - Relación romántica compleja
  - Se conocieron durante el periodo de aprendizaje
  - Tensión ideológica: ella quiere revolución violenta, él quiere reforma
  - Lo ama pero teme que el sistema lo corrompa
  - Conflicto central: ¿Puede el amor sobrevivir a ideologías opuestas?
\begin{itemize}
  \item \textbf{Rol en la historia}: 
\end{itemize}
  - Interés romántico principal
  - Voz de la resistencia radical
  - Ayuda a Ethan con información y ajustes técnicos de El Fénix
  - Representa la resistencia activa vs. el conformismo
  - Tensión dramática cuando Ethan gana El Velo
\begin{itemize}
  \item \textbf{Habilidades}: Explosivos, sabotaje, conocimiento del mercado negro

  \item \textbf{Frase característica}: "Todo lo que vuela puede caer"

  \item \textbf{Arco narrativo}: De antagonista ideológica a aliada romántica conflictuada
\end{itemize}

#### EL TOCAYO (CARLOS)
\begin{itemize}
  \item \textbf{Edad}: Viejo (~70 años)

  \item \textbf{Profesión}: Comerciante en la Plaza de Chatarra

  \item \textbf{Rol}: Proveedor de piezas, contacto en el mercado negro

  \item \textbf{Personalidad}: Astuto, oportunista, pero con código de honor
\end{itemize}

#### VIPER
\begin{itemize}
  \item \textbf{Edad}: ~25-30 años

  \item \textbf{Rol}: Piloto corporativo, antagonista principal en la carrera

  \item \textbf{Descripción}: Campeón patrocinado, arrogante, despiadado

  \item \textbf{Vehículo}: Motowave de última generación corporativa

  \item \textbf{Personalidad}: Elitista, ve a los del Distrito RD como basura
\end{itemize}

---

\subsection}{FAMILIA (AUSENTE/MISTERIO)}

#### MADRE DE ETHAN
\begin{itemize}
  \item \textbf{Nombre}: [A definir]
  \item \textbf{Estado Actual}: Operando en los niveles más profundos de Goglee City bajo el más estricto secreto.
  \item \textbf{Historia}: Fue reclutada a la fuerza (o mediante un trato desesperado) por el Alto Mando corporativo en 2100, tras descubrir ella misma las perturbaciones del Agujero Negro Errante.
  \item \textbf{Rol Verdadero}: Directora técnica del Proyecto Noé. Está diseñando las Arcas Gravitacionales.
  \item \textbf{Relación con Ethan}: Orquestadora en la sombra. Ella filtra tecnología y desafíos específicos al Distrito RD (a través del Tocayo o canales anónimos) para "entrenar" a Ethan sin que él lo sepa. Ella sabe que Ethan es el único con la sensibilidad neuronal para pilotar el Arca principal en la tormenta final.
  \item \textbf{Misterio}: Su supuesta desaparición fue la tapadera perfecta para su reclutamiento corporativo.

\end{itemize}

---

\section{EL SECRETO DE LA CÚPULA: PROYECTO NOÉ Y EL CATACLISMO ESTELAR}

\subsection{EL AGUJERO NEGRO ERRANTE (LA SOMBRA VIAJERA)}
\begin{itemize}
  \item \textbf{La Verdadera Amenaza}: Los astrónomos de la cúpula han detectado señales confusas de perturbaciones gravitacionales en sistemas estelares lejanos. Los datos indican la presencia de un \textbf{Agujero Negro Errante} moviéndose a una velocidad extrema hacia el Sistema Solar.
  \item \textbf{Escala del Desastre}: No es solo un aumento del nivel del mar; es la destrucción potencial de la Tierra, el Sol y quizás todo el sistema. La Gran Inundación de 2058 fue solo un preámbulo provocado por las primeras ondas de marea gravitacional.
  \item \textbf{El Ocultamiento}: Las megacorporaciones mantienen esta "Fecha Final" bajo el más estricto secreto para evitar el colapso social absoluto mientras ellos construyen en secreto los "Arcas" (el verdadero propósito detrás de la tecnología de El Velo).
\end{itemize}

\subsection{EL PROYECTO NOÉ}
\begin{itemize}
  \item \textbf{La Misión de la Madre}: La madre de Ethan descubrió la data real y entendió que la élite solo planeaba salvarse a sí misma. El "Proyecto Noé" es su plan para asegurar que el ADN humano y la cultura del "fango" (la verdadera resiliencia) sobrevivan.
  \item \textbf{Ethan como el Noé Moderno}: Cada prueba de pilotaje, cada geodésica surfeada, es una preparación para que Ethan pueda pilotar una de las Arcas a través de la tormenta gravitacional del Agujero Negro.
  \item \textbf{El Factor Humano Indispensable}: Las Arcas del Proyecto Noé utilizan motores de curvatura que generan agujeros de gusano artificiales. Las IAs y computadoras de navegación colapsan ante las fluctuaciones de la singularidad; solo un piloto capaz de "ver" y "sentir" las geodésicas en tiempo real puede corregir el rumbo y evitar la espaguetización del Arca. Ethan es ese piloto.
  \item \textbf{Sincronicidad Mística}: Los sueños premonitorios de Carlos y la "conexión" de Ethan con las geodésicas son en realidad la sensibilidad de sus sistemas nerviosos a las perturbaciones del espacio-tiempo que se aproxima.
\end{itemize}

#### FAMILIA EXTENDIDA RAMÍREZ
\begin{itemize}
  \item \textbf{Padre/Madre de Yeisly}: Hermano/hermana de Carlos

  \item \textbf{Otros primos}: Posible comunidad familiar más amplia en Distrito RD

  \item \textbf{Abuelos de Ethan}: Refugiados originales de República Dominicana (fallecidos)
\end{itemize}

---

\subsection\textit{{PERSONAJES DEL PASADO (FALLECIDOS)}

#### LUNA REYES (2060-2084) †
\begin{itemize}
  \item \textbf{Edad al morir}: 24 años

  \item \textbf{Fecha de nacimiento}: ~2060

  \item \textbf{Fecha de muerte}: 2084 (durante El Velo Gravitacional)

  \item \textbf{Descripción física}: [A desarrollar en flashbacks]

  \item \textbf{Personalidad}: 
\end{itemize}
  - Piloto prodigio
  - Valiente hasta la temeridad
  - Soñadora y optimista
  - Creía que podía cambiar el mundo
  - Veía el espacio-tiempo de forma diferente a los demás
\begin{itemize}
  \item \textbf{Relación con Tomás}: Mejor amiga y compañera de carrera

  \item \textbf{Habilidad especial}: Única piloto que dominaba El Surfeo de Geodésicas
\end{itemize}
  - Podía "ver" la geometría del espacio-tiempo
  - Identificaba geodésicas que otros no podían percibir
  - Ganó 3 carreras usando rutas "imposibles"
\begin{itemize}
  \item \textbf{Muerte}: 
\end{itemize}
  - Murió en la carrera de El Velo 2084
  - Intentó surfear una geodésica extremadamente peligrosa a través de un cañón
  - Error de cálculo de 3 milímetros en el ángulo de entrada
  - \textbf{Efecto Honda} la expulsó a 800 km/h hacia atrás
  - Colisionó con un Martillo Gravitacional
  - Murió gritando mientras su cuerpo era aplastado por fuerzas G extremas
  - Tomás escuchó sus gritos por radio durante 30 segundos
  - Vio su motowave desintegrarse en la pantalla
\begin{itemize}
  \item \textbf{Legado}: 
\end{itemize}
  - Su técnica se convirtió en leyenda prohibida
  - Goglee Corporation prohibió oficialmente El Surfeo de Geodésicas
  - Nadie más intentó la técnica después de su muerte (considerada "suicidio elegante")
  - Hasta Ethan Ramírez en 2114
\begin{itemize}
  \item \textbf{Frases características}: 
\end{itemize}
  - "Es como escuchar una canción que nadie más puede oír"
  - "Ves valles y crestas donde otros solo ven aire vacío"
  - "El espacio-tiempo no es plano, es un océano de olas invisibles"
\begin{itemize}
  \item \textbf{Impacto en la historia}: 
\end{itemize}
  - Su muerte traumatizó a Tomás profundamente
  - Tomás dejó de competir después de perderla
  - Hace que enseñar la técnica a Ethan sea un dilema moral devastador
  - Ethan representa una segunda oportunidad para Tomás

---

\section}{LUGARES Y GEOGRAFÍA}

\subsection\textit{{GOGLEE CITY}
\begin{itemize}
  \item \textbf{Tipo}: Megaciudad bajo domo protector

  \item \textbf{Fundación}: 2048

  \item \textbf{Ubicación}: [Continente a definir]

  \item \textbf{Corporación}: Goglee Corporation

  \item \textbf{Características}:
\end{itemize}
  - Domo de protección ambiental
  - Aire limpio y filtrado
  - Luz artificial controlada
  - Temperatura regulada
  - Arquitectura futurista
\begin{itemize}
  \item \textbf{Población}: Ciudadanos privilegiados con "Ciudadanía Aérea"

  \item \textbf{Gobierno}: Corporatocracia

  \item \textbf{Economía}: Tecnología avanzada, servicios de lujo

  \item \textbf{Sociedad}: Estratificada, elitista
\end{itemize}

\subsection}{DISTRITO RD}
\begin{itemize}
  \item \textbf{Nombre completo}: Distrito República Dominicana

  \item \textbf{Ubicación}: Periferia EXTERIOR del domo de Goglee City

  \item \textbf{Fundación}: 2060 (refugiados de La Gran Inundación)

  \item \textbf{Características}:
\end{itemize}
  - Fuera del domo protector
  - Aire tóxico pero respirable
  - Sol visible (enfermo, rojizo)
  - Smog permanente
  - Infraestructura deteriorada
  - Neón parpadeante
  - Charcos de agua ácida
\begin{itemize}
  \item \textbf{Población}: 
\end{itemize}
  - Descendientes de refugiados caribeños (principalmente dominicanos)
  - Segunda y tercera generación
  - Mantienen cultura y tradiciones de la isla
\begin{itemize}
  \item \textbf{Economía}: 
\end{itemize}
  - Vertedero industrial
  - Mano de obra para Goglee City
  - Mercado negro
\begin{itemize}
  \item \textbf{Lugares importantes}:
\end{itemize}
  - \textbf{Plaza de Chatarra}: Mercado de piezas y desechos
  - \textbf{Taller de Carlos}: Pequeño agujero en edificio condenado

\subsection\textit{{REGIONES CARIBEÑAS SUMERGIDAS (2058)}
\begin{itemize}
  \item \textbf{Estado}: Sumergidas o devastadas por La Gran Inundación.

  \item \textbf{Islas del Caribe}: Cuba, República Dominicana, Haití, Puerto Rico, Jamaica, Bahamas (todas bajo el agua).

  \item \textbf{Miami}: Ahora es un arrecife de coral artificial bajo cien metros de agua salada.

  \item \textbf{Venezuela}: Ciudades costeras convertidas en infraestructuras náufragas; migración masiva a los Andes y al Distrito RD.

  \item \textbf{Colombia}: Playas devoradas por un mar que no pidió perdón; Cartagena y Barranquilla son leyendas bajo las olas.

  \item \textbf{Dáspora}: La mezcla de sangres dominicana, cubana, venezolana y colombiana define la cultura del Distrito RD.
\end{itemize}

---

\section}{TECNOLOGÍA}

\subsection\textit{{MOTOWAVES}
\begin{itemize}
  \item \textbf{Definición}: Vehículos híbridos rueda/levitación

  \item \textbf{Funcionamiento}:
\end{itemize}
  - Modo terrestre: Ruedas convencionales
  - Modo levitación: Propulsores gravitacionales
\begin{itemize}
  \item \textbf{Desarrollo}: Comercializados desde ~2075

  \item \textbf{Tipos}:
\end{itemize}
  - Corporativos: Alta tecnología, nuevos, letales
  - Artesanales: Modificados a mano, únicos

\subsection}{EL FÉNIX}
\begin{itemize}
  \item \textbf{Tipo}: Motowave modificada artesanalmente

  \item \textbf{Constructor}: Ethan Ramírez (con guía de Carlos)

  \item \textbf{Historia}: 
\end{itemize}
  - Base: Motowave usada (regalo de cumpleaños de Carlos)
  - Transformación: Febrero-diciembre 2114
\begin{itemize}
  \item \textbf{Apariencia}: Estéticamente hermosa (NO chatarra)

  \item \textbf{Características}:
\end{itemize}
  - Ruedas gastadas para asfalto roto
  - Propulsores gravitacionales (luz azul)
  - Estabilizador de fase T-9 (grado militar)
\begin{itemize}
  \item \textbf{Modificaciones}:
\end{itemize}
  - Sistema de levitación artesanal
  - Estabilizador de fase integrado
  - Ajustes personalizados para el estilo de Ethan
\begin{itemize}
  \item \textbf{Sonido característico}: Ronroneo suave, zumbido grave al levitar
\end{itemize}

\subsection\textit{{ESTABILIZADOR DE FASE}
\begin{itemize}
  \item \textbf{Modelo}: T-9 (grado militar)

  \item \textbf{Función (Estructural/Pasiva)}: 
\end{itemize}
  - Previene la desintegración molecular durante la inversión gravitacional extrema.
  - Sincroniza la frecuencia de la motowave con la geodésica para evitar que el vehículo "atraviese" la pista.
  - Es un elemento de \textbf{supervivencia estática}.
\begin{itemize}
  \item \textbf{Importancia}: Crítico para sobrevivir El Velo; sin él, la moto se desintegraría bajo fuerzas de marea.

  \item \textbf{Costo}: 500-1000 créditos (fortuna para Distrito RD)
\end{itemize}

\subsection}{INVERSOR G}
\begin{itemize}
  \item \textbf{Función (Activa/Maniobra)}: 
\end{itemize}
  - Permite al piloto invertir la polaridad de la gravedad local instantáneamente.
  - Actúa como el "volante" gravitatorio para maniobras suicidas.
  - Permite pegarse a paredes, realizar giros bruscos o frenado de emergencia mediante repulsión.
\begin{itemize}
  \item \textbf{Propuesto por}: Unika, como mejora para compensar la falta de tecnología corporativa.

  \item \textbf{Riesgo}: Un retraso en la activación (latencia) puede lanzar al piloto contra el suelo antes de que la gravedad cambie.
\end{itemize}

\subsection\textit{{PROPULSORES GRAVITACIONALES}
\begin{itemize}
  \item \textbf{Función}: Levitación vehicular

  \item \textbf{Apariencia}: Luz azul tenue

  \item \textbf{Desarrollo}: 2065-2075

  \item \textbf{Mecanismo}: [A desarrollar técnicamente]
\end{itemize}

\subsection}{MARTILLOS GRAVITACIONALES}
\begin{itemize}
  \item \textbf{Uso}: Obstáculos en El Velo Gravitacional

  \item \textbf{Función}: Pistones gravitacionales invisibles que aplastan

  \item \textbf{Ubicación}: Sector 4 de la pista

  \item \textbf{Estrategia}: Ir por los bordes donde la gravedad es más débil
\end{itemize}

\subsection\textit{{DOMOS PROTECTORES}
\begin{itemize}
  \item \textbf{Función}: Protección ambiental de ciudades.

  \item \textbf{Aisalmiento}: El aire interior es filtrado y "estéril", lo que quema los pulmones de quienes están acostumbrados al Distrito.

  \item \textbf{Goglee City}: El domo principal, rodeado de muros de cristal reforzado y sensores tácticos.
\end{itemize}

---

\subsection}{TÉCNICAS GRAVITACIONALES AVANZADAS}

#### EL SURFEO DE GEODÉSICAS (Geodesic Surfing) / LA LÍNEA FANTASMA

\begin{itemize}
  \item \textbf{Clasificación}: Técnica avanzada / Leyenda

  \item \textbf{Nivel de peligro}: EXTREMO (Mortalidad 90%+)

  \item \textbf{Base científica}: Relatividad General (geodésicas en espacio-tiempo curvado)

  \item \textbf{Descubrimiento}: ~2080 (por accidente durante experimentos gravitacionales)

  \item \textbf{Única maestra conocida}: Luna Reyes (2060-2084) †
\end{itemize}

\textbf{Fundamento físico:}

En la Relatividad General de Einstein, el camino más corto entre dos puntos en un espacio curvado se llama \textbf{geodésica}. Aunque para un observador externo parezca una curva larga, para quien la recorre es la línea más recta y eficiente posible.

Los círculos flotantes de El Velo no son solo metas, sino \textbf{generadores de micro-curvatura} que alteran el tejido del espacio-tiempo entre puntos.

\textbf{Descripción técnica:}

1. \textbf{Identificar geodésicas ocultas}: Detectar caminos curvos en el espacio-tiempo que son más cortos que las líneas rectas aparentes
2. \textbf{Caer en pozos gravitatorios}: Usar las micro-curvaturas generadas por los círculos para obtener aceleración "gratuita"
3. \textbf{Comprimir distancia efectiva}: En términos relativistas, recorrer una distancia menor porque el espacio está comprimido
4. \textbf{Surfear la curvatura}: Deslizarse por el espacio-tiempo curvado como un surfista en una ola

\textbf{Mecánica paso a paso:}

1. \textbf{Escaneo}: El Lidar-G mapea las fluctuaciones de densidad gravitacional en tiempo real
2. \textbf{Cálculo}: El piloto identifica la geodésica óptima (requiere intuición espacial extrema)
3. \textbf{Entrada}: Alinear la motowave con el ángulo exacto de la geodésica
4. \textbf{Surfeo}: Dejarse "caer" en el pozo gravitatorio, ganando velocidad sin combustible
5. \textbf{Salida}: Emerger en el punto objetivo con velocidad y posición ventajosas

\textbf{Tecnología requerida:}

| Componente | Función |
|------------|---------|
| \textbf{Escáner Lidar-G} | Mapea fluctuaciones de densidad gravitacional en tiempo real. Muestra la "topografía" invisible del espacio-tiempo |
| \textbf{Compensador de Inercia} | Permite tomar curvas de 90° a velocidades supersónicas sin que el cuerpo sea aplastado por fuerzas G extremas |
| \textbf{Velas de Espacio-Tiempo} | Pequeños alerones electromagnéticos que "muerden" la curvatura para maniobrar sin necesidad de aire |
| \textbf{Núcleo Gravitacional} | Genera campo de protección contra distorsiones extremas |

\textbf{Cómo se siente:}

\begin{itemize}
  \item Luna Reyes: \textit{"Es como escuchar una canción que nadie más puede oír. Ves valles y crestas donde otros solo ven aire vacío."}

  \item Tomás: \textit{"El universo entero se vuelve visible. Puedes ver cómo el espacio se dobla, se retuerce, te llama."}

  \item Requiere la habilidad de \textbf{"sentir" la geometría del espacio-tiempo} instintivamente
\end{itemize}

\textbf{Peligros mortales:}

1. \textbf{Efecto Honda (Slingshot Effect)}
   - Error de cálculo por milímetros en el ángulo de entrada
   - La curvatura te "escupe" hacia afuera a velocidad inmanejable
   - Similar al efecto honda gravitacional usado en viajes espaciales, pero invertido
   - Resultado: Colisión fatal o expulsión al vacío

2. \textbf{Dilatación Temporal Local}
   - En curvaturas extremas, el tiempo fluye diferente
   - Puedes ganar en tu reloj pero llegar tarde en el reloj de la meta
   - Resultado: Descalificación o peor

3. \textbf{Colapso Estructural}
   - Fuerzas G extremas en la transición geodésica
   - La motowave se desintegra si no está reforzada adecuadamente
   - Resultado: Muerte instantánea

4. \textbf{Efecto de Lente Gravitacional}
   - Desapareces de la vista por el doblamiento de la luz
   - Otros pilotos no te ven y pueden colisionar
   - Resultado: Accidente multi-piloto

\textbf{Ventaja en la carrera:}

\begin{itemize}
  \item Acortar distancias efectivas en 30-40%

  \item Aceleración "gratuita" por caída gravitatoria

  \item Aparecer en lugares "imposibles" instantáneamente

  \item Ventaja de 15-20 segundos (diferencia entre ganar y perder)

  \item Parecer "invisible" por efecto de lente gravitacional
\end{itemize}

\textbf{Historia:}

\begin{itemize}
  \item \textbf{2080}: Técnica descubierta por accidente durante experimentos

  \item \textbf{2082-2084}: Luna Reyes domina la técnica
\end{itemize}
  - Usó la técnica exitosamente 12 veces en carreras
  - Ganó 3 carreras usando geodésicas imposibles
  - Se convirtió en leyenda del Distrito RD
\begin{itemize}
  \item \textbf{2084}: Muerte de Luna Reyes
\end{itemize}
  - Error de cálculo de 3 milímetros en ángulo de entrada
  - Efecto Honda la expulsó a 800 km/h hacia atrás
  - Colisionó con un Martillo Gravitacional
  - Murió gritando mientras su cuerpo era aplastado
  - Tomás escuchó todo por radio durante 30 segundos
  - Dejó de competir después de su muerte
\begin{itemize}
  \item \textbf{2084-2114}: Nadie más intentó la técnica
\end{itemize}
  - Considerada "suicidio elegante"
  - Goglee Corporation la prohibió oficialmente
  - Se convirtió en leyenda urbana
\begin{itemize}
  \item \textbf{2114}: Ethan Ramírez aprende la técnica
\end{itemize}
  - Tomás decide enseñarle después de ver su don natural
  - Ethan puede "ver" las geodésicas instintivamente

\textbf{Estado actual:}

\begin{itemize}
  \item Técnica prohibida por Goglee Corporation

  \item Considerada leyenda urbana por pilotos jóvenes

  \item Solo Tomás "El Viejo Toro" Vega conoce los detalles completos

  \item Debe decidir si enseñársela a Ethan
\end{itemize}

\textbf{Dilema moral de Tomás:}

\begin{itemize}
  \item Sin El Surfeo de Geodésicas, Ethan no puede ganar contra pilotos corporativos con tecnología superior

  \item Con El Surfeo de Geodésicas, Ethan puede morir como Luna

  \item Tomás debe elegir entre proteger a Ethan o darle la única herramienta para triunfar

  \item El trauma de escuchar morir a Luna vs. la esperanza de que Ethan lo logre
\end{itemize}

\textbf{Escena icónica (ejemplo narrativo):}

\textit{Ethan apaga la guía de navegación oficial. Ve dos círculos que parecen estar en extremos opuestos de un cañón. Los otros pilotos rodean el cañón por la línea "recta" aparente, perdiendo segundos preciosos.}

\textit{Ethan, en cambio, ve algo que ellos no pueden ver: una geodésica latente, un valle invisible en el espacio-tiempo que conecta ambos puntos directamente a través del cañón.}

\textit{Se lanza hacia lo que parece un vacío. Por un segundo, desaparece de la vista debido al efecto de lente gravitacional. El espacio se dobla a su alrededor.}

\textit{Y aparece de repente del otro lado, cruzando el círculo antes que nadie, como si hubiera teletransportado.}

---

#### OTRAS TÉCNICAS (Menos peligrosas)

\textbf{Sobrecarga de Núcleo:}
\begin{itemize}
  \item Desactivar limitadores de seguridad temporalmente

  \item Aumenta velocidad 200-300%

  \item Riesgo: Explosión del motor

  \item Ethan usa esta técnica en el Capítulo 3
\end{itemize}

\textbf{Deslizamiento Gravitacional:}
\begin{itemize}
  \item Usar campos gravitacionales como "rampas"

  \item Permite giros imposibles

  \item Técnica estándar de pilotos avanzados
\end{itemize}

\textbf{Inversión de Polaridad:}
\begin{itemize}
  \item Cambiar la dirección de los propulsores instantáneamente

  \item Permite frenado de emergencia

  \item Riesgo: Daño estructural a la motowave
\end{itemize}

---

\section\textit{{CULTURA Y TRADICIONES}

\subsection}{BENDICIÓN TRADICIONAL CARIBEÑA}
\begin{itemize}
  \item \textbf{Origen}: Isla del Caribe (pre-inundación)

  \item \textbf{Cion Papi:} Un intercambio verbal de respeto y bendición arraigado en la cultura caribeña. El hijo solicita la bendición ("Cion, papi") y el padre responde con un deseo de protección ("Dios te bendiga, mijo"). Es un vínculo de reconocimiento mutuo que no requiere de parafernalia física ni religiosa formal, aunque puede ir acompañado de un gesto natural de afecto como un abrazo o un apretón de hombros.

  \item \textbf{Naturaleza}: Ritual verbal/gestual de respeto (NO es literal besar la mano)

  \item \textbf{Ritual}:
\end{itemize}
  1. Descendiente se acerca al ascendiente (padre, tío, tía, abuelo, etc.)
  2. Descendiente dice: \textbf{"La bendición, papi/tío/tía/abuela"}
  3. Forma abreviada común: \textbf{"Ción, papi"} / \textbf{"Ción, papa"} / \textbf{"Ción, tío"}
  4. Ascendiente responde: \textbf{"Dios te bendiga"} / \textbf{"Dios te bendiga, hijo/mijo"} / \textbf{"Dios me lo bendiga"}
\begin{itemize}
  \item \textbf{Contexto}: 
\end{itemize}
  - Carlos es ateo pero mantiene la tradición cultural
  - Representa conexión con raíces caribeñas y respeto familiar
  - Se usa en momentos importantes (despedidas, antes de eventos peligrosos)
  - Tradición específica de la cultura de la isla y del Distrito RD

\subsection\textit{{IDENTIDAD CULTURAL DEL DISTRITO RD}
\begin{itemize}
  \item \textbf{Herencia}: Caribeña mezclada (refugiados de islas y costas continentales)
\end{itemize}
  - Mezcla de culturas: dominicana, cubana, venezolana, colombiana, puertorriqueña, haitiana, etc.
  - Imposible distinguir orígenes específicos después de 3 generaciones
  - El Distrito RD es un crisol cultural caribeño
\begin{itemize}
  \item \textbf{Idioma}: Español mezclado con inglés

  \item \textbf{Slang caribeño} (usado por jóvenes):
\end{itemize}
  - \textbf{"Klk"}, \textbf{"Klk manito"} o \textbf{"Klk manín"}: Saludo informal y versátil. Es la marca de identidad de Rafa.
\begin{itemize}
  \item \textbf{"Manín" / "Manito"}: Término de hermandad (diminutivo de hermano).

  \item \textbf{"Hey! Klk"}: Saludo casual común entre la juventud del Distrito RD.
\end{itemize}
  - Mezcla natural de español, inglés y expresiones caribeñas de diversas islas
\begin{itemize}
  \item \textbf{Valores}: Familia, comunidad, resistencia

  \item \textbf{Memoria colectiva}: Historias de la isla perdida

  \item \textbf{Música}: [A desarrollar]

  \item \textbf{Comida}: [A desarrollar]
\end{itemize}

---

\section}{HISTORIA Y CRONOLOGÍA}

\subsection\textit{{LÍNEA DE TIEMPO GLOBAL}

#### \textbf{2025-2045: INDUSTRIALIZACIÓN EXTREMA}
\begin{itemize}
  \item Aceleración del cambio climático antropogénico

  \item Smog permanente cubre ciudades principales

  \item Temperatura global aumenta +3°C

  \item Corporaciones ganan poder sobre gobiernos
\end{itemize}

#### \textbf{2040-2055: ERA DE LOS DOMOS}
\begin{itemize}
  \item \textbf{2040}: Primeros prototipos de domos corporativos

  \item \textbf{2045}: Primer domo funcional (ciudad pequeña)

  \item \textbf{2048}: \textbf{Fundación de Goglee City} (domo corporativo)

  \item \textbf{2045-2055}: Construcción de domos en continentes principales
\end{itemize}

#### \textbf{2055-2060: LA GRAN INUNDACIÓN}
\begin{itemize}
  \item \textbf{2052}: Temperatura global +3°C sobre niveles pre-industriales

  \item \textbf{2055}: Inicio del derretimiento acelerado del Ártico

  \item \textbf{2058}: \textbf{LA GRAN INUNDACIÓN}
\end{itemize}
  - Colapso completo del casquete polar ártico
  - Nivel del mar sube 7-10 metros
  - \textbf{Islas del Caribe completamente sumergidas}: Cuba, República Dominicana, Haití, Puerto Rico, Jamaica, Bahamas
  - \textbf{Costas continentales devastadas}: Miami y sur de Florida, costa caribeña de Venezuela y Colombia, Centroamérica
  - ~45 millones de refugiados solo de islas caribeñas (muchos más de costas continentales; muchos murieron)
\begin{itemize}
  \item \textbf{2060}: Fundación oficial del \textbf{Distrito RD} (periferia de Goglee City)
\end{itemize}

#### \textbf{2060-2080: POST-INUNDACIÓN}
\begin{itemize}
  \item \textbf{2060-2070}: Asentamiento de refugiados en periferias

  \item \textbf{2065}: Desarrollo de tecnología de levitación vehicular

  \item \textbf{2070}: Primeros prototipos de propulsores gravitacionales

  \item \textbf{2075}: Comercialización limitada de motowaves

  \item \textbf{2078}: Goglee Corp lanza \textbf{El Velo Gravitacional} (primera edición)
\end{itemize}

#### \textbf{2080-2114: ERA ACTUAL}
\begin{itemize}
  \item Consolidación de la sociedad de domos

  \item Distrito RD se convierte en ghetto permanente

  \item El Velo Gravitacional como espectáculo anual
\end{itemize}

---

\subsection}{LÍNEA DE TIEMPO FAMILIAR RAMÍREZ}

```
2058 (3 de enero)    - Nacimiento de CARLOS RAMÍREZ
                       (mismo año de La Gran Inundación)
                       Hijo de refugiados dominicanos

2058-2093            - Carlos crece en Distrito RD
                     - Aprende mecánica
                     - Se convierte en maestro

~2093                - Carlos conoce a la madre de Ethan
                       (Carlos tiene ~35 años)

2095 (3 de febrero)  - Nacimiento de ETHAN RAMÍREZ
                       (Carlos tiene 37 años y 1 mes)

~2100                - Desaparición de la madre de Ethan
                       (Ethan tiene ~5 años, Carlos ~42 años)

2100-2114            - Carlos cría solo a Ethan
                     - Le enseña mecánica y cultura caribeña

2114 (3 de enero)    - Cumpleaños de Carlos (56 años)

2114 (3 de febrero)  - Cumpleaños de Ethan (17 años)
                     - Carlos le regala una motowave usada

2114 (feb-dic)       - PERIODO DE APRENDIZAJE
                     - Ethan transforma la moto en "El Fénix"
                     - 10 meses de construcción y entrenamiento

2114 (diciembre)     - EVENTOS DE LA HISTORIA PRINCIPAL
                     - Ethan compite en El Velo Gravitacional
                     - Gana la carrera
                     - Rechaza la ciudadanía
                     - Inicia la revolución
```

---

\section\textit{{EVENTOS CLAVE}

\subsection}{LA GRAN INUNDACIÓN (2058)}
\begin{itemize}
  \item \textbf{Causa}: Derretimiento completo del casquete polar ártico

  \item \textbf{Consecuencia inmediata}: Nivel del mar +7-10 metros

  \item \textbf{Regiones afectadas}: 
\end{itemize}
  - \textbf{Islas del Caribe}: Completamente sumergidas (Cuba, República Dominicana, Haití, Puerto Rico, Jamaica, Bahamas, islas menores)
  - \textbf{Costas continentales}: Miami y sur de Florida, costa caribeña de Venezuela, costa caribeña de Colombia, costa este de Centroamérica, zonas del Golfo de México
  - \textbf{Costas globales}: Inundaciones masivas en todos los continentes
\begin{itemize}
  \item \textbf{Víctimas}: Cientos de millones globalmente

  \item \textbf{Refugiados}: ~45 millones solo del Caribe (muchos más de costas continentales)

  \item \textbf{Impacto social}: Creación de ghettos de refugiados (como Distrito RD)
\end{itemize}

\subsection\textit{{EL VELO GRAVITACIONAL (CARRERA)}
\begin{itemize}
  \item \textbf{Primera edición}: 2078

  \item \textbf{Frecuencia}: Anual

  \item \textbf{Organizador}: Goglee Corporation

  \item \textbf{Propósito oficial}: Espectáculo deportivo

  \item \textbf{Propósito real}: 
\end{itemize}
  - Entretenimiento para élite
  - Válvula de escape para marginados (falsa esperanza)
  - Propaganda corporativa
\begin{itemize}
  \item \textbf{Premio}: Ciudadanía Aérea (acceso al domo)

  \item \textbf{Peligros}:
\end{itemize}
  - Inversión gravitacional extrema
  - Martillos Gravitacionales (Sector 4)
  - Sabotaje entre competidores
  - Alta tasa de mortalidad
\begin{itemize}
  \item \textbf{Pista}: Serpentea hacia las puertas del domo de Goglee City
\end{itemize}

\subsection}{EL REGALO (3 de febrero 2114)}
\begin{itemize}
  \item Carlos regala a Ethan una motowave usada

  \item Momento emotivo padre-hijo

  \item Inicio del viaje de Ethan hacia El Velo
\end{itemize}

\subsection\textit{{LA VICTORIA Y EL RECHAZO (Diciembre 2114)}
\begin{itemize}
  \item Ethan gana El Velo Gravitacional

  \item Rechaza la Ciudadanía Aérea

  \item Declaración pública contra el sistema

  \item Inicio de movimiento revolucionario
\end{itemize}

---

\section}{TEMAS Y MOTIVOS}

\subsection\textit{{TEMAS PRINCIPALES}

1. \textbf{Desigualdad Social y Estratificación}
   - Domo vs. Periferia
   - Ciudadanos vs. Refugiados
   - Acceso a recursos y tecnología

2. \textbf{Identidad Cultural y Memoria}
   - Preservación de tradiciones caribeñas
   - Generaciones que nunca vieron la isla
   - Tensión entre asimilación y resistencia

3. \textbf{Cambio Climático y Responsabilidad}
   - Consecuencias del capitalismo descontrolado
   - Corporaciones que causan el problema y venden la "solución"
   - Quiénes pagan el precio del "progreso"

4. \textbf{Familia y Legado}
   - Relación padre-hijo (Carlos-Ethan)
   - Transmisión de conocimiento y valores
   - Ausencia de la madre

5. \textbf{Resistencia vs. Conformismo}
   - Ethan (reforma desde dentro) vs. Sony (revolución desde fuera)
   - El costo de la esperanza
   - Cuándo es necesario romper el sistema

6. \textbf{Tecnología como Doble Filo}
   - Avances que salvan y oprimen
   - Acceso desigual a la tecnología
   - Creatividad artesanal vs. producción corporativa

\subsection}{MOTIVOS RECURRENTES}

\begin{itemize}
  \item \textbf{Gravedad}: Metáfora de opresión y liberación

  \item \textbf{El Fénix}: Renacimiento, transformación

  \item \textbf{Luz azul}: Esperanza, tecnología, lo artificial

  \item \textbf{El sol enfermo}: Mundo natural destruido

  \item \textbf{Bendición tradicional}: Conexión con raíces, amor familiar

  \item \textbf{Volar}: Libertad, aspiración, peligro
\end{itemize}

---

\section\textit{{NOTAS DE DESARROLLO}

\subsection}{ELEMENTOS A DESARROLLAR}

\begin{itemize}
  \item [ ] Descripción física detallada de Ethan

  \item [ ] Descripción física de Yeisly, Rafa, Unica, Marcos

  \item [ ] Nombre y backstory completa de la madre de Ethan

  \item [ ] Mecanismo técnico de propulsores gravitacionales

  \item [ ] Mecanismo técnico de estabilizadores de fase

  \item [ ] Música y comida del Distrito RD

  \item [x] Personajes del círculo cercano de Ethan (Yeisly, Rafa, Tomás, Unika, Marcos)

  \item [ ] Detalles de la pista de El Velo

  \item [ ] Estructura política de Goglee City

  \item [ ] Otras ciudades-domo en el mundo

  \item [ ] Desarrollo de la revolución post-rechazo

  \item [ ] Historia completa de Tomás "El Viejo Toro" Vega (su victoria en 2082, castigo, pérdidas)

  \item [ ] Familia extendida Ramírez (padres de Yeisly, otros primos)

  \item [ ] Subplot del arte de Yeisly como resistencia

  \item [ ] Desarrollo de la posible tensión romántica Ethan-Unica (admiración secreta de ella)
\end{itemize}

\subsection*{PREGUNTAS ABIERTAS}

\begin{itemize}
  \item ¿En qué continente está Goglee City?

  \item ¿Qué pasó exactamente con la madre de Ethan?

  \item ¿Cuántos competidores hay en El Velo?

  \item ¿Qué otras corporaciones existen?

  \item ¿Hay resistencia organizada antes de Ethan?
\end{itemize}

---

\textbf{FIN DE LA BIBLIA DEL UNIVERSO v1.3}

\textbf{Changelog v1.3}:
\begin{itemize}
  \item \textbf{CAMBIO MAYOR}: Reemplazada técnica "El Anclaje Fantasma" con \textbf{"El Surfeo de Geodésicas"}

  \item Basada en física real: Relatividad General y geodésicas en espacio-tiempo curvado

  \item Agregada tecnología: Escáner Lidar-G, Compensador de Inercia, Velas de Espacio-Tiempo

  \item Actualizada muerte de Luna Reyes: Efecto Honda por error de 3mm (más científicamente plausible)

  \item Agregados nuevos peligros: Efecto Honda, Dilatación Temporal, Colapso Estructural, Efecto de Lente

  \item Incluida escena icónica de ejemplo (cañón geodésico)

  \item Actualizado perfil de Tomás y Luna con nueva mecánica
\end{itemize}

\textbf{Changelog v1.2}:
\begin{itemize}
  \item Agregada técnica \textbf{El Anclaje Fantasma} (Ghost Anchor) con detalles completos

  \item Agregado personaje \textbf{Luna Reyes} (piloto fallecida, 2060-2084)

  \item Expandida historia de Tomás "El Viejo Toro" Vega con tragedia de Luna

  \item Agregada sección de Técnicas Gravitacionales Avanzadas

  \item Actualizada edad de Sony a 19-20 años (novia de Ethan)

  \item Cambiada dinámica: Unica admira secretamente a Ethan desde pequeña (él no lo sabe)
\end{itemize}

\textbf{Changelog v1.1}:
\begin{itemize}
  \item Agregados 5 nuevos personajes secundarios del Periodo de Aprendizaje

  \item Expandida sección de familia extendida Ramírez

  \item Agregados detalles de Tomás "El Viejo Toro" Vega como precedente histórico
\end{itemize}

\end{document}
